% !TEX root = ../paper.tex
\section{GPU Pipeline}
\label{sec:pipeline}

The GPU implementation provides two execution paths selected at runtime based
on the current shell particle count $N$. Both paths share the capture stage and
the host-side core update; they differ in how shell physics is dispatched.

\subsection{Large-$N$ Path ($N > 2048$)}
\label{sec:large_n}

Seven stages execute on two CUDA streams with fork/join synchronisation via
lightweight events (\texttt{cudaEventDisableTiming}):

\begin{table}[H]
  \caption{Large-$N$ GPU pipeline stages.}
  \label{tab:pipeline_large}
  \centering
  \scriptsize
  \begin{tabular}{@{}clccc@{}}
    \toprule
    \textbf{\#} & \textbf{Stage} & \textbf{Kernel} & \textbf{Stream} & \textbf{Block} \\
    \midrule
    1 & Capture         & \texttt{captureFromSnowpack}  & sim  & 256 \\
    2 & Core update     & host code                     & ---  & --- \\
    3 & Forces+Tether   & \texttt{forceTether}          & sim  & 256 \\
    4 & Grid build      & hash$\to$sort$\to$ranges      & grid & 256 \\
    5 & Collision       & \texttt{neighborCollision}    & sim  & 512 \\
    6 & Integrate+Ground& \texttt{integrateGround}      & sim  & 256 \\
    7 & Core RB         & host code                     & ---  & --- \\
    \bottomrule
  \end{tabular}
\end{table}

Stages 3 (forces\,+\,tether) and 4 (grid build) execute \textbf{concurrently}
on \texttt{simStream} and \texttt{gridStream}, respectively. The collision
kernel (stage~5) waits for both via \texttt{cudaStreamWaitEvent} on
\texttt{evGridDone}. Stream synchronisation events use
\texttt{cudaEventDisableTiming} to eliminate timing overhead (Section~\ref{sec:streams}).

\subsection{Small-$N$ Path ($N \leq 2048$)}
\label{sec:small_n}

When the shell count is small, the overhead of the grid build (5--8 GPU
commands: clear + hash + CUB sort + cell ranges) exceeds the computational
benefit. A single \textbf{mega-fused kernel} (\texttt{shellPhysicsFused})
replaces stages 3--6, reducing GPU commands per frame from $\sim$10 to
\textbf{1} and keeping all intermediate forces in registers
(Section~\ref{sec:mega_fused}).

\subsection{Block Sizes and \texttt{\_\_launch\_bounds\_\_}}
\label{sec:block_sizes}

Block sizes are chosen to balance \textbf{SM coverage} (all SMs receive at
least one block) against per-block resource consumption. With $N = 4{,}000$ and
5~SMs on the GTX~960M, a block of 1024 yields only $\lceil4000/1024\rceil = 4$
blocks, leaving 1~SM idle (20\% wasted). A block of 256 yields 16~blocks,
ensuring full coverage.

The \texttt{neighborCollisionKernel} uses block size 512 dictated by the
warp-cooperative tile layout ($\text{MAX\_WARPS} = 16 = 512/32$). All other
physics kernels use 256. Each kernel is annotated with
\texttt{\_\_launch\_bounds\_\_(maxThreads, minBlocks)} to guide the register
allocator (Table~\ref{tab:occupancy}).

\subsection{Adaptive Path Selection}
\label{sec:adaptive_path}

The $N = 2048$ threshold is chosen empirically: below this value, the
mega-fused kernel's reduced command overhead ($\sim$100--150\,$\mu$s) outweighs
the cost of brute-force $O(N^2)$ collision detection. Above 2048, the grid-based
$O(N \cdot k)$ path becomes more efficient despite the additional GPU commands.
At typical steady-state conditions ($N \approx 4{,}000$ particles), the grid path
executes $\sim$3--4 neighbours per particle, yielding $\sim$16{,}000 pair
evaluations compared to $\sim 16\text{M}$ in brute-force --- a $1000\times$
reduction that justifies the grid-build overhead.
